%! Carrying Out a Test for a Population Mean — Unit 7, Lesson 7.5 — Warm-Up (Answer Key)
\documentclass[10pt]{article}
\usepackage{apstats-article}
\usepackage{apstats-key}   % pulls in -boxes; do NOT also load -boxes
\newcommand{\TallMath}[1]{$\displaystyle #1\rule[-1.4em]{0pt}{3.2em}$}

\begin{document}

\pageheader{Unit 7, Lesson 7.5}{Warm-Up --- Answer Key}
\namedateperiod

\vspace{0.15in}

\begin{notesbox}{Spiral Review --- Setting Up a $t$-Test (Lesson 7.4)}

A nutritionist claims the mean daily sodium intake of adults is 2,400 mg. A health researcher
believes the true mean is \emph{higher}. She takes a random sample of 35 adults.
$\bar x = 2{,}510$ mg and $s = 380$ mg. A histogram of the data shows a roughly symmetric
distribution with no outliers.

\textbf{1.} Name the inference method and give the degrees of freedom.

\quad Method: \ans{one-sample $t$-test for $\mu$} \quad $df = $ \ans{34}

\textbf{2.} State hypotheses.

\quad $H_0$: \ans{$\mu = 2400$} \quad $H_a$: \ans{$\mu > 2400$}

\textbf{3.} Verify both conditions.
\begin{itemize}
  \item Random: \ansline{Random sample of 35 adults --- condition met.}
  \item Normal: \ansline{$n = 35 \ge 30$; CLT applies --- condition met.}
\end{itemize}

\textbf{4.} Using the formula $t = \dfrac{\bar x - \mu_0}{s / \sqrt{n}}$, compute the
$t$-statistic. Show your work.

\vspace{0.15in}
$SE = \dfrac{s}{\sqrt{n}} = \dfrac{{\color{keyred}\mathbf{380}}}{\sqrt{{\color{keyred}\mathbf{35}}}} \approx $ \ans{64.24}
\hfill
$t = \dfrac{{\color{keyred}\mathbf{2510}} - {\color{keyred}\mathbf{2400}}}{{\color{keyred}\mathbf{64.24}}} \approx $ \ans{1.71}

\end{notesbox}

\end{document}
